% Generated draft: requires mathematical and source review.
\documentclass[11pt]{article}
\usepackage[margin=25mm]{geometry}
\usepackage{fontspec}
\setmainfont{Latin Modern Roman}
\usepackage{amsmath,amssymb,mathtools,unicode-math}
\setmathfont{Latin Modern Math}
\usepackage{enumitem,needspace,xurl,hyperref}
\setlength{\parindent}{0pt}
\setlength{\parskip}{5pt}
\setlength{\emergencystretch}{3em}
\newcommand{\N}{\mathbb{N}}
\newcommand{\Z}{\mathbb{Z}}
\newcommand{\Q}{\mathbb{Q}}
\newcommand{\R}{\mathbb{R}}
\newcommand{\C}{\mathbb{C}}
\newcommand{\F}{\mathbb{F}}
\newcommand{\A}{\mathbb{A}}
\newcommand{\PP}{\mathbb{P}}
\newcommand{\E}{\mathbb{E}}
\newcommand{\eps}{\varepsilon}
\newcommand{\dd}{\,\mathrm{d}}
\begin{document}
\section*{UnsolvedMath ID: 30006543}
\textbf{Source problem code:} catalogue entry
\subsection{Definitions and mathematical statement}
Let \(P=(X,\leq_P)\) be a finite partially ordered set, let \(n=|X|\geq 1\), and let \(x\in X\). The supplied record asks whether the source-specific statistic \(\operatorname{win}_P(x)\) and the rank variance \(\sigma_P^2(x)\) satisfy
\[
\operatorname{win}_P(x)=\Omega\!\left(\frac{\sigma_P^2(x)}{n}\right)
\]
uniformly over \(P\) and \(x\). This is identified in the supplied report excerpt as Conjecture 10.

The record does not define \(\operatorname{win}_P(x)\), so a complete mathematical formulation cannot be recovered from the supplied text. It also does not explicitly specify the probability distribution used for rank variance. Under the usual interpretation, a linear extension is a bijection \(L:X\to\{1,\ldots,n\}\) such that \(u<_P v\) implies \(L(u)<L(v)\). If \(L\) is selected uniformly from the finite set of linear extensions of \(P\), the rank of \(x\) is \(R_x=L(x)\), and its variance is
\[
\sigma_P^2(x)=\mathbb{E}\!\left[\left(R_x-\mathbb{E}[R_x]\right)^2\right].
\]
This probability model is an explicitly stated possible interpretation requiring confirmation from the source, not a definition supplied by the record.

With that interpretation and the source's definition of \(\operatorname{win}\) restored, a conventional uniform asymptotic reading of the conjecture would be
\[
\exists c>0\;\exists n_0\in\N\;\forall n\geq n_0\;
\forall P\text{ with }|P|=n\;\forall x\in P,
\qquad
\operatorname{win}_P(x)\geq c\,\frac{\sigma_P^2(x)}{n}.
\]
Here \(c\) and \(n_0\) are independent of \(P\), \(x\), and \(n\). The source's wording ``for every \(P\) and \(x\)'' may instead intend the bound for every \(n\geq1\), with a universal \(c>0\). The supplied record does not resolve this distinction.

The excerpt further reports that even the following asymptotic possibility was not known to be excluded: there are posets \(P_k\), elements \(x_k\in P_k\), sizes \(n_k=|P_k|\to\infty\), and constants \(a,B>0\) such that, for all sufficiently large \(k\),
\[
\sigma_{P_k}^2(x_k)\geq a n_k^2,
\qquad
\operatorname{win}_{P_k}(x_k)\leq B.
\]
The excerpt describes these as respectively the largest and smallest possible orders of the two quantities. These are supplied contextual claims, not independently verified bounds.

The surrounding excerpt also describes qualitative local implications involving width, \(\operatorname{win}\), rank standard deviation, and a parameter \(\pi\): a large antichain must contain an element with large \(\operatorname{win}\), and large \(\operatorname{win}_P(x)\) implies large \(\sigma_P(x)\). It describes an implication involving \(\pi\) and \(\sigma\) as nonlocal, giving the example of two unrelated chains of size \(n/2\), with \(n\) even and \(x\) minimal in one chain. The implication symbols, quantitative meanings, and definition of \(\pi\) are missing. Subsequent comparisons among Theorems 3--6 and the displayed conditions in Theorem 6 are also incomplete or corrupted, including an undefined conclusion numbered (3). They cannot be reconstructed as precise auxiliary statements from this record.
\subsection{Short English statement}
Conjecture 10 asks whether every element of every finite poset has its source-defined win statistic bounded below by a universal constant times its rank variance divided by the poset size. The definition of the win statistic, the rank probability model, and the exact asymptotic quantifiers require confirmation. The supplied dated curation labels the conjecture open; this draft does not independently verify that status.
\subsection{Sources}
\begin{itemize}
\item Supplied local UnsolvedMath registry record, permanent ID \(30006543\), title: Linear-Extension Wins and Poset Variance; exact source problem code: \texttt{catalogue entry}. Collection: Oberwolfach Reports Open Problems. The record assigns proposed year 2026 and supplies no individual proposer.
\item Supplied primary source citation: \emph{Combinatorics} (2026), Oberwolfach Reports, DOI 10.4171/owr/2026/1, pp. 5--79. The supplied excerpt identifies the question as Conjecture 10. \url{https://doi.org/10.4171/owr/2026/1}.
\item Supplied literature source: Aires--Kahn, \emph{Balancing Extensions in Posets of Large Width}. \url{https://arxiv.org/abs/2509.11549}.
\item Literature assessment dated exactly \texttt{2026-08-22}: the supplied curation labels Conjecture 10 open and states that available results establish implications among width, window, and variance parameters without proving the pointwise bound \(\operatorname{win}_P(x)=\Omega(\sigma_P^2(x)/n)\). This is supplied dated curation, not an independently verified literature assessment. The record itself requests independent verification of time-sensitive status and cited claims.
\item The registry supplies the statement-status label \texttt{corrected\_verified} and reports that the official Oberwolfach report was checked and the notation repaired. These are supplied provenance claims; no source visit or independent mathematical verification is asserted here.
\end{itemize}
\end{document}
